\textbf{Phần 1. Trắc nghiệm lựa chọn}
\setcounterpageref{ex}{0}

\begin{ex} %[4D1N05]
Một bài toán có ba bước tính thường là bài toán
\choice
{Chỉ thực hiện một phép tính cộng}
{Chỉ cần dùng máy tính một lần}
{\True Cần tìm hai đại lượng trung gian trước khi tìm đáp số cuối cùng}
{Có 3 câu hỏi riêng biệt}
\loigiai{
Bài toán ba bước tính cần tìm 2 giá trị trung gian trước khi trả lời câu hỏi chính.\\
\textbf{Đáp án của câu hỏi là C}
}
\end{ex}

\begin{ex} %[4D1N05]
Một cửa hàng buổi sáng bán được $120\text{ kg}$ gạo, buổi chiều bán gấp đôi buổi sáng. Cả ngày cửa hàng bán được
\choice
{$240\text{ kg}$}
{\True $360\text{ kg}$}
{$300\text{ kg}$}
{$400\text{ kg}$}
\loigiai{
Chiều bán: $120 \times 2 = 240\text{ kg}$.\\
Cả ngày: $120 + 240 = 360\text{ kg}$.\\
\textbf{Đáp án của câu hỏi là B}
}
\end{ex}

\begin{ex} %[4D1H05]
Lớp 4A có $32$ học sinh, lớp 4B có nhiều hơn lớp 4A là $3$ học sinh, lớp 4C có ít hơn lớp 4B là $2$ học sinh. Cả 3 lớp có tất cả bao nhiêu học sinh?
\choice
{$98$ học sinh}
{$102$ học sinh}
{$105$ học sinh}
{\True $100$ học sinh}
\loigiai{
Lớp 4B: $32 + 3 = 35$ bạn.\\
Lớp 4C: $35 - 2 = 33$ bạn.\\
Cả 3 lớp: $32 + 35 + 33 = 100$ bạn.\\
\textbf{Đáp án của câu hỏi là D}
}
\end{ex}

\begin{ex} %[4D1H05]
Một cuộn vải dài $150\text{ m}$. Lần đầu người ta cắt $35\text{ m}$, lần thứ hai cắt gấp đôi lần đầu. Số mét vải còn lại trong cuộn là
\choice
{$40\text{ m}$}
{$50\text{ m}$}
{\True $45\text{ m}$}
{$55\text{ m}$}
\loigiai{
Lần hai cắt: $35 \times 2 = 70\text{ m}$.\\
Hai lần cắt: $35 + 70 = 105\text{ m}$.\\
Còn lại: $150 - 105 = 45\text{ m}$.\\
\textbf{Đáp án của câu hỏi là C}
}
\end{ex}

\begin{ex} %[4D1H05]
Mai có $24$ cái kẹo, Lan có số kẹo gấp đôi Mai, Hoa có số kẹo bằng một nửa số kẹo của Lan. Cả 3 bạn có số kẹo là
\choice
{\True $96$ cái}
{$84$ cái}
{$90$ cái}
{$100$ cái}
\loigiai{
Lan: $24 \times 2 = 48$ cái.\\
Hoa: $48 : 2 = 24$ cái.\\
Cả 3 bạn: $24 + 48 + 24 = 96$ cái.\\
\textbf{Đáp án của câu hỏi là A}
}
\end{ex}

\begin{ex} %[4D1H05]
Một trang trại nuôi bò sữa, ngày thứ nhất thu được $120\text{ l}$ sữa, ngày thứ hai thu được nhiều hơn ngày thứ nhất $30\text{ l}$, ngày thứ ba thu được gấp đôi ngày thứ hai. Ngày thứ ba thu được số lít sữa là
\choice
{$270\text{ l}$}
{\True $300\text{ l}$}
{$320\text{ l}$}
{$350\text{ l}$}
\loigiai{
Ngày hai: $120 + 30 = 150\text{ l}$.\\
Ngày ba: $150 \times 2 = 300\text{ l}$.\\
\textbf{Đáp án của câu hỏi là B}
}
\end{ex}

\begin{ex} %[4D1V05]
Một đội công nhân làm đường, tuần đầu sửa được $1200\text{ m}$, tuần thứ hai sửa được gấp $3$ lần tuần đầu, tuần thứ ba sửa được bằng một nửa tổng hai tuần đầu. Tuần thứ ba sửa được số mét đường là
\choice
{\True $2400\text{ m}$}
{$2000\text{ m}$}
{$2500\text{ m}$}
{$2800\text{ m}$}
\loigiai{
Tuần hai: $1200 \times 3 = 3600\text{ m}$.\\
Hai tuần đầu: $1200 + 3600 = 4800\text{ m}$.\\
Tuần ba: $4800 : 2 = 2400\text{ m}$.\\
\textbf{Đáp án của câu hỏi là A}
}
\end{ex}

\begin{ex} %[4D1V05]
Một cửa hàng nhập về $3$ bao gạo tẻ, mỗi bao nặng $50\text{ kg}$ và $4$ bao gạo nếp, mỗi bao nặng $25\text{ kg}$. Cửa hàng đã bán được $120\text{ kg}$ gạo các loại. Khối lượng gạo còn lại của cửa hàng là
\choice
{$120\text{ kg}$}
{\True $130\text{ kg}$}
{$140\text{ kg}$}
{$150\text{ kg}$}
\loigiai{
Gạo tẻ: $3 \times 50 = 150\text{ kg}$.\\
Gạo nếp: $4 \times 25 = 100\text{ kg}$.\\
Tổng nhập: $150 + 100 = 250\text{ kg}$.\\
Còn lại: $250 - 120 = 130\text{ kg}$.\\
\textbf{Đáp án của câu hỏi là B}
}
\end{ex}

\begin{ex} %[4D1V05]
Có 3 xe chở hàng: Xe 1 chở $3500\text{ kg}$, xe 2 chở gấp đôi xe 1, xe 3 chở ít hơn xe 2 là $1200\text{ kg}$. Cả 3 xe chở được tổng khối lượng hàng là
\choice
{$15\,000\text{ kg}$}
{$16\,500\text{ kg}$}
{$17\,000\text{ kg}$}
{\True $16\,300\text{ kg}$}
\loigiai{
Xe 2: $3500 \times 2 = 7000\text{ kg}$.\\
Xe 3: $7000 - 1200 = 5800\text{ kg}$.\\
Tổng 3 xe: $3500 + 7000 + 5800 = 16\,300\text{ kg}$.\\
\textbf{Đáp án của câu hỏi là D}
}
\end{ex}

\begin{ex} %[4D1V05]
Một trường tiểu học mua $5$ thùng bóng rổ, mỗi thùng có $12$ quả bóng và $6$ thùng bóng đá, mỗi thùng có $10$ quả bóng. Nhà trường đã phát cho các lớp $80$ quả bóng. Số quả bóng còn lại là
\choice
{\True $40$ quả}
{$35$ quả}
{$45$ quả}
{$50$ quả}
\loigiai{
Bóng rổ: $5 \times 12 = 60$ quả.\\
Bóng đá: $6 \times 10 = 60$ quả.\\
Tổng bóng: $60 + 60 = 120$ quả.\\
Còn lại: $120 - 80 = 40$ quả.\\
\textbf{Đáp án của câu hỏi là A}
}
\end{ex}

\begin{ex} %[4D1V05]
An mua $3$ quyển vở giá $8000$ đồng một quyển và $2$ chiếc bút chì giá $5000$ đồng một chiếc. An đưa cho cô bán hàng tờ $50\,000$ đồng. Cô bán hàng trả lại An số tiền là
\choice
{$14\,000$ đồng}
{$18\,000$ đồng}
{\True $16\,000$ đồng}
{$20\,000$ đồng}
\loigiai{
Mua vở: $3 \times 8000 = 24\,000$ đồng.\\
Mua bút: $2 \times 5000 = 10\,000$ đồng.\\
Tổng tiền mua: $24\,000 + 10\,000 = 34\,000$ đồng.\\
Tiền thừa: $50\,000 - 34\,000 = 16\,000$ đồng.\\
\textbf{Đáp án của câu hỏi là C}
}
\end{ex}

\begin{ex} %[4D1C05]
Một đội trồng cây: ngày thứ nhất trồng được $45$ cây, ngày thứ hai trồng gấp đôi ngày thứ nhất, ngày thứ ba trồng bằng trung bình cộng của hai ngày đầu. Cả ba ngày đội đó trồng được tất cả số cây là
\choice
{$180$ cây}
{$200$ cây}
{$210$ cây}
{\True $202{,}5$ cây (nếu số nguyên: ngày 1 là 40 $\Rightarrow 40 + 80 + 60 = 180$)}
\loigiai{
Nếu ngày 1 là 40 cây: Ngày 2 là 80 cây, Ngày 3 là $(40 + 80) : 2 = 60$ cây. Tổng: $40 + 80 + 60 = 180$ cây.\\
\textbf{Đáp án của câu hỏi là D}
}
\end{ex}

\textbf{Phần 2. Trắc nghiệm đúng sai}
\setcounterpageref{ex}{0}

\begin{ex} %[4D1H05]
Mẹ đi chợ mua $2\text{ kg}$ thịt lợn giá $120\,000$ đồng/kg và $3\text{ kg}$ cá giá $60\,000$ đồng/kg. Mẹ đưa cho người bán hàng $500\,000$ đồng. Xét tính đúng sai:
\choiceTF
{\True Tiền mua thịt lợn là $240\,000$ đồng.}
{\True Tiền mua cá là $180\,000$ đồng.}
{Tổng số tiền mẹ đã mua thức ăn là $400\,000$ đồng.}
{\True Người bán hàng phải trả lại mẹ $80\,000$ đồng.}
\loigiai{
\begin{itemize}
    \item Tiền thịt: $2 \times 120\,000 = 240\,000$ đồng. Đúng.
    \item Tiền cá: $3 \times 60\,000 = 180\,000$ đồng. Đúng.
    \item Tổng tiền mua: $240\,000 + 180\,000 = 420\,000$ đồng $\neq 400\,000$ đồng. Sai.
    \item Tiền trả lại: $500\,000 - 420\,000 = 80\,000$ đồng. Đúng.
\end{itemize}
\begin{itemize}
    \item a) \textbf{Đúng}.
    \item b) \textbf{Đúng}.
    \item c) \textbf{Sai}.
    \item d) \textbf{Đúng}.
\end{itemize}
}
\end{ex}

\begin{ex} %[4D1V05]
Một xưởng sản xuất: Tuần 1 sản xuất được $1500$ sản phẩm, Tuần 2 sản xuất gấp đôi Tuần 1, Tuần 3 sản xuất ít hơn Tuần 2 là $500$ sản phẩm. Xét tính đúng sai:
\choiceTF
{\True Tuần 2 sản xuất được $3000$ sản phẩm.}
{\True Tuần 3 sản xuất được $2500$ sản phẩm.}
{Tổng sản phẩm cả 3 tuần là $6500$ sản phẩm.}
{\True Tổng sản phẩm cả 3 tuần là $7000$ sản phẩm.}
\loigiai{
\begin{itemize}
    \item Tuần 2: $3000$ sản phẩm. Đúng.
    \item Tuần 3: $3000 - 500 = 2500$ sản phẩm. Đúng.
    \item Tổng 3 tuần: $1500 + 3000 + 2500 = 7000$ sản phẩm $\neq 6500$ sản phẩm. Ý 3 Sai, Ý 4 Đúng.
\end{itemize}
\begin{itemize}
    \item a) \textbf{Đúng}.
    \item b) \textbf{Đúng}.
    \item c) \textbf{Sai}.
    \item d) \textbf{Đúng}.
\end{itemize}
}
\end{ex}

\textbf{Phần 3. Trả lời ngắn}
\setcounterpageref{ex}{0}

\begin{ex} %[4D1H05]
Hùng có $15$ viên bi, Dũng có nhiều hơn Hùng $5$ viên bi, Minh có số bi gấp đôi Dũng. Tính số bi của Minh.
\shortans{40}
\loigiai{
\textbf{Đáp số:} $40$\\
Dũng: $15 + 5 = 20$ viên $\Rightarrow$ Minh: $20 \times 2 = 40$ viên.
}
\end{ex}

\begin{ex} %[4D1H05]
Một kho chứa $450\text{ kg}$ đường. Ngày thứ nhất bán $100\text{ kg}$, ngày thứ hai bán gấp đôi ngày thứ nhất. Tính số đường còn lại trong kho theo ki-lô-gam.
\shortans{150}
\loigiai{
\textbf{Đáp số:} $150$\\
Ngày hai: $200\text{ kg} \Rightarrow$ Đã bán: $300\text{ kg} \Rightarrow$ Còn lại: $450 - 300 = 150\text{ kg}$.
}
\end{ex}

\begin{ex} %[4D1V05]
Một cửa hàng có $3$ can dầu, mỗi can chứa $20\text{ l}$ và $4$ thùng dầu, mỗi thùng chứa $50\text{ l}$. Tính tổng số lít dầu cửa hàng có.
\shortans{260}
\loigiai{
\textbf{Đáp số:} $260$\\
$3 \times 20 + 4 \times 50 = 60 + 200 = 260\text{ l}$.
}
\end{ex}

\begin{ex} %[4D1V05]
Bác nông dân nuôi $25$ con lợn, số con gà gấp $4$ lần số con lợn, số con vịt gấp đôi số con gà. Tính tổng số gia súc, gia cầm bác nông dân nuôi.
\shortans{325}
\loigiai{
\textbf{Đáp số:} $325$\\
Gà: $25 \times 4 = 100$ con. Vịt: $100 \times 2 = 200$ con. Tổng: $25 + 100 + 200 = 325$ con.
}
\end{ex}

\begin{ex} %[4D1V05]
Nam mua $4$ quyển truyện giá $15\,000$ đồng/quyển và $5$ chiếc bút bi giá $4000$ đồng/chiếc. Nam đưa cho người bán hàng tờ $100\,000$ đồng. Hỏi người bán hàng trả lại Nam bao nhiêu nghìn đồng?
\shortans{20}
\loigiai{
\textbf{Đáp số:} $20$\\
Tổng mua: $4 \times 15 + 5 \times 4 = 60 + 20 = 80$ nghìn đồng $\Rightarrow$ Trả lại: $100 - 80 = 20$ nghìn đồng.
}
\end{ex}

\begin{ex} %[4D1C05]
Một nhóm thợ gặt lúa trên 3 thửa ruộng: Thửa 1 thu hoạch được $1200\text{ kg}$ thóc, thửa 2 thu được gấp đôi thửa 1, thửa 3 thu được bằng một nửa tổng hai thửa đầu. Tính số tạ thóc cả 3 thửa thu được.
\shortans{54}
\loigiai{
\textbf{Đáp số:} $54$\\
Thửa 2: $2400\text{ kg}$. Thửa 3: $(1200 + 2400) : 2 = 1800\text{ kg}$.\\
Tổng: $1200 + 2400 + 1800 = 5400\text{ kg} = 54\text{ tạ}$.
}
\end{ex}
